\documentclass{ceurart}

\usepackage{listings}
\usepackage{tabularx}
\newcommand{\term}[1]{\texttt{#1}}
\lstset{basicstyle=\small\ttfamily,breaklines=true,columns=fullflexible,
  keepspaces=true,showstringspaces=false,frame=single,aboveskip=6pt,belowskip=6pt}
\emergencystretch=1em
\sloppy

\begin{document}

\copyrightyear{2027}
\copyrightclause{Copyright for this paper by its authors.
  Use permitted under Creative Commons License Attribution 4.0
  International (CC BY 4.0).}
\conference{SWAT4HCLS 2027, February 1--4, 2027, Basel, Switzerland}

\title{VCF Core: A Semantic Foundation for RDF Representations of VCF}

\author[1,2]{Elias Crum}[
orcid=0009-0005-3991-754X,
email=elias.crum@ugent.be
]
\cormark[1]
\author[2]{Bart Buelens}[
orcid=0000-0001-7734-3747,
email=bart.buelens@vito.be
]
\author[2]{Gökhan Ertaylan}[
orcid=0000-0001-5602-6435,
email=gokhan.ertaylan@vito.be
]
\author[1]{Ruben Taelman}[
orcid=0000-0001-5118-256X,
email=ruben.taelman@ugent.be
]
\address[1]{IDLab, Department of Electronics and Information Systems, Ghent University -- imec, Belgium}
\address[2]{Flemish institute for Technological Research (VITO) Mol, Belgium}
\cortext[1]{Corresponding author.}


\begin{abstract}
Medical uses of genomic data increasingly depend on integrating variant observations with clinical information and biomedical knowledge.
Semantic data linking offers a promising basis for this integration, but the widely used Variant Call Format (VCF) does not natively expose its contents as Linked Data, and existing RDF renderings are either partial or tied to one converter.
We present the VCF Core Vocabulary, a converter-independent semantic target for VCF, and the design decisions behind it.

Three commitments shape the model.
It retains the source context of an observation, so that file metadata, header declarations, records, annotations and genotype data remain distinguishable and interpretable.
It reuses published vocabularies instead of re-minting their concepts: FALDO supplies positions and strandedness outright, while the Sequence Ontology, GA4GH VRS, GENO, ChEBI and HERO are linked as alignment targets rather than asserted as replacements for VCF resources.
And it treats version differences as constraint differences rather than model differences, so one vocabulary serves VCF 4.1 through 4.5 while five SHACL conformance profiles carry the per-version rules.
Two complementary genotype profiles remain available: an expanded representation exposing individual sample calls, and a condensed representation grouping them into sample-ordered vectors.

We report coverage against VCF 4.5 from a published construct inventory rather than an impression: 100 of 106 logical-model constructs are fully covered, four partially, and two are absent, alongside all 122 reserved INFO and FORMAT key declarations.
Scope boundaries---BCF byte layout, byte-for-byte reconstruction---are stated explicitly rather than left implicit.
Every statistic in this paper is recomputed from the artifacts by a script in the repository, which fails if the vocabulary and the inventory diverge.
\end{abstract}


\begin{keywords}
Variant Call Format\sep RDF\sep Genomic data\sep Vocabulary\sep SHACL
\end{keywords}
\maketitle

\section{Introduction and related work}

The Variant Call Format (VCF) is the de-facto interchange format for genomic variant data~\cite{danecek2011vcf}. A VCF record is interpreted through more than its position and alleles: header declarations describe annotations, the FORMAT column determines the contents of sample fields, and file metadata supplies reference and processing context~\cite{vcf45}. A reusable RDF representation must therefore address both the variation described and the source observation from which that description derives.

Two workloads make genotype representation especially consequential. For a single-sample VCF, explicit genotype and read-depth statements support focused inspection and integration. In a cohort VCF, the same representation creates resources across the product of records, samples and FORMAT fields. A vocabulary should make this choice explicit, so that different conversion systems can share a semantic target while selecting an appropriate degree of materialization.

Existing work provides complementary foundations. GFVO harmonizes GFF3, GTF, GVF and VCF concepts in one ontology and already renders VCF contents in RDF~\cite{baran2015gfvo}, where we instead retain what one file declares. VCF2RDF describes an isomorphic mapping from VCF to RDF~\cite{penha2017vcf2rdf}. HERO-Genomics integrates genomic entities with wider biomedical information and models VCF files and records~\cite{cazzaro2025hero}. Semantic Beacons combines ontology-based mappings with federated queries over genomic variation and public knowledge graphs~\cite{bodrug2025semantic}. The separately maintained gvar schema describes genomic variants as Linked Data and explicitly distinguishes that purpose from representing VCF itself~\cite{gvar}. The Sequence Ontology (SO) and FALDO supply reusable concepts for sequence features and their locations~\cite{eilbeck2005so,bolleman2016faldo}, and the GA4GH Variation Representation Specification (VRS) supplies a computable model of variation with federated identifiers~\cite{wagner2021vrs}. Our earlier conversion-framework paper proposed an ontology for single-sample VCF transformation~\cite{crum2025semantifying}.

This paper develops that direction through the VCF Core Vocabulary 2.0.0~\cite{vcfc200}, renamed from the VCF-RDFizer Vocabulary so that a shared semantic target is not identified by one converter's name; the retired \term{vcfr:} namespace resolves to a deprecation document in which every 1.1.0 term remains defined, deprecated and linked to its successor. We contribute a VCF-facing model that retains source context and header-linked interpretation (Section~\ref{sec:design}); a deliberate division of labour with published vocabularies, reusing FALDO outright and aligning to SO, VRS, GENO, ChEBI and HERO rather than restating them (Section~\ref{sec:reuse}); two complementary genotype profiles, \term{vcfc:ExpandedRepresentation} and \term{vcfc:CondensedRepresentation} (Section~\ref{sec:profiles}); and a measured, reproducible account of how much of VCF 4.5 the model covers, with SHACL conformance profiles that let one vocabulary serve VCF 4.1 through 4.5 (Section~\ref{sec:coverage}). Throughout, \term{vcfc:} denotes \url{https://w3id.org/vcf-core/vocab#}.

\section{Scope and design rationale}
\label{sec:design}

\subsection{Retaining the source context}

The vocabulary separates a \term{VCFFile}, its \term{VCFHeader} and header lines, individual \term{VCFRecord} resources, and the assessment associated with each record, represented by a \term{VariantCall}. The latter groups QUAL, FILTER, INFO, FORMAT and genotype data; it is a modeling abstraction rather than an additional VCF line type. Figure~\ref{fig:model} summarizes the shared structure and the genotype alternatives.

This separation preserves distinctions between observations from different files. Records describing an equivalent alteration can carry different quality scores, filters, annotations or sample calls, and collapsing them into a single normalized variant resource would discard exactly the evidence a downstream consumer may need to weigh. Accordingly, the file is modeled as a \term{dcat:Distribution} and a \term{prov:Entity}; records, calls and sample resources are provenance entities as well, providing attachment points for provenance descriptions~\cite{prov2013}. The recommended file-based IRI templates retain the connection between source, record, call and field. They identify source-local resources rather than globally normalized biological variants. Producers must choose unambiguous record identifiers: neither position alone nor a potentially missing VCF ID guarantees row identity.

The property \term{asSequenceAlteration} links a record or call to an independently modeled SO sequence alteration (\term{SO:0001059}). This allows biological interpretation, and alignment with models such as gvar, to evolve without replacing the source representation. Such links require appropriate reference and allele interpretation; their presence does not establish automatic cross-model equivalence. A \term{VCFSample} similarly identifies a sample column within a file, without assuming that the column is itself a biological individual---a distinction that matters for tumour/normal pairs and for pooled or synthetic columns.

\subsection{Reusing published vocabularies rather than restating them}
\label{sec:reuse}

The vocabulary follows one rule, recorded in its own ontology header: a term is minted only when it carries VCF syntax that has no equivalent in a published vocabulary. Everything else is delegated. The delegation happens at two deliberately distinct strengths, and we keep them visibly separate because they make different commitments and carry different risks.

The stronger form is reuse inside the vocabulary's own axioms, where an external term constrains what VCF Core means. Five external terms appear in \term{rdfs:subClassOf}, \term{rdfs:subPropertyOf} or \term{rdfs:range} position. A \term{VCFFile} is a \term{dcat:Distribution}; files, records, calls, sample calls, sample sets, cohort matrices and value vectors are \term{prov:Entity} instances; the clonal \term{Original=} relation of a \term{\#\#PEDIGREE} line is a subproperty of \term{prov:wasDerivedFrom}, which is precisely what that line asserts; \term{asSequenceAlteration} ranges over \term{SO:0001059}; and the confidence intervals around POS and END range over \term{faldo:InRangePosition}.

FALDO is reused more strongly still, being used directly in data rather than only referenced. Record and allele positions are \term{faldo:location} links to \term{faldo:ExactPosition} or \term{faldo:Region} resources whose \term{faldo:reference} names the declared contig. Imprecise structural-variant bounds---the CIPOS and CIEND confidence intervals---are \term{faldo:InRangePosition} resources with \term{faldo:begin} and \term{faldo:end}, which is exactly FALDO's idiom for a position known only within a range. The per-strand value pairs of the VCF 4.5 base-modification fields use \term{faldo:ForwardStrandPosition} and \term{faldo:ReverseStrandPosition}, and the specification's unstranded convention maps onto \term{faldo:BothStrandsPosition}. The consequence is concrete and independently checkable: the vocabulary declares no position, region or strand class of its own, and a consumer that already understands FALDO can locate every VCF Core record without learning a second location model.

The weaker form is alignment. Forty-five external terms occur only as objects of \term{skos:exactMatch}, \term{skos:closeMatch}, \term{skos:relatedMatch} or \term{rdfs:seeAlso}: SO~(15), ChEBI~(10), VRS~(8), HERO-Genomics~(6), FALDO~(3) and GENO~(3). Several of these matches are close enough to be worth naming. VRS \term{Adjacency} corresponds to a VCF breakend and \term{Terminus} to a single or telomeric breakend; \term{CisPhasedBlock} corresponds to a phase set; and \term{DerivativeMolecule} with \term{TraversalBlock} corresponds to the derivative-chromosome traversal that VCF 4.5 encodes in the PSO field. The nine symbolic structural-variant allele types map onto SO classes such as \term{SO:0000159} (deletion) and \term{SO:1000173} (tandem duplication). The ten base-modification alias keys carry the ChEBI identifiers that the specification itself names, so \term{M5mC} aligns with \term{CHEBI:27551}. GENO supplies close matches for the genotype layer, and HERO-Genomics for the file and record layer, where an existing model already covers the same ground.

Alignment is deliberately not assertion. A VCF record is an observation made by one caller on one file; the same alteration observed elsewhere may carry different quality, filters and annotations. Asserting VRS or SO class membership on the VCF resource would make the source representation depend on a biological interpretation that the file does not itself justify, and would silently commit every consumer of the graph to that interpretation. Separating the two strengths also means an alignment can be corrected without disturbing the VCF model. This matters in practice: an earlier release aligned contig declarations to a \term{faldo:Reference} class that does not exist, and correcting it touched one annotation rather than the model. Where no honest match exists we leave the term unaligned---the four reserved EVENTTYPE codes \term{DUP:DISPERSED}, \term{TRA:UNBALANCED}, \term{BFB} and \term{DOUBLEMINUTE} have no SO counterpart and are not forced onto a near-match.

\begin{figure}[t]
\centering\small
\fbox{\parbox{0.92\linewidth}{\centering
\textbf{Shared VCF context}\\[1pt]
\term{VCFFile} $\rightarrow$ \term{VCFRecord} $\rightarrow$ \term{VariantCall};
\term{VCFHeader} contains header lines and field definitions}}
\par\smallskip
\begin{tabularx}{\linewidth}{@{}>{\centering\arraybackslash}X >{\centering\arraybackslash}X@{}}
$\downarrow$ \textbf{Expanded} & $\downarrow$ \textbf{Condensed} \\
\term{SampleCall} & \term{CohortCallMatrix} \\
$\downarrow$ \term{hasFormatValue} & $\downarrow$ \term{hasFormatValueVector} \\
\term{FormatFieldValue} & \term{FormatValueVector} \\
One resource per sample/key & One vector per key, ordered by \term{SampleSet}
\end{tabularx}
\caption{Selected conceptual structure. Both genotype branches attach to a \term{VariantCall}. Expanded values can link to FORMAT definitions through \term{declaredBy}; condensed vectors must. The file declares the chosen profile.}
\label{fig:model}
\end{figure}

\subsection{Header-defined meaning and controlled interpretation}

Header information is represented explicitly because it determines how values should be interpreted. All eleven meta-information line types defined by VCF 4.5 have dedicated classes---fileformat, INFO, FILTER, FORMAT, ALT, assembly, contig, META, SAMPLE, PEDIGREE and pedigreeDB---as does the \term{\#CHROM} column line. Because the specification also permits structured lines it does not define, a generic \term{HeaderAttribute} carries the key-value attributes of any structured line, so producer-specific and future declarations remain queryable without a bespoke predicate for each. INFO and FORMAT definitions expose \term{fieldId}, \term{fieldNumber}, \term{fieldType} and \term{fieldDescription}, together with the recommended \term{fieldSource} and \term{fieldVersion}; structured values link to their definitions through \term{declaredBy}, and condensed vectors use the same relation.

One distinction is worth drawing out because it is easy to get wrong. A reserved key such as INFO/DP is a specification-level declaration, whereas a \term{\#\#INFO} line is a fact about one file. \term{FieldDefinition} is therefore not a subclass of \term{HeaderLine}; \term{declaredOnLine} links a definition to the header line that declares it in a given file. Conflating the two makes every entry of a shipped reserved-key registry appear to be a header line that is missing its key and attributes.

Reusing a generic field-value pattern avoids introducing a vocabulary predicate for every producer-specific annotation, and it keeps identically named keys distinct when their declaration or scope differs: INFO/DP and FORMAT/DP need not describe the same measurement. Retaining the local definition supports interpretation, although harmonizing meanings across producers still requires additional mappings. \term{fieldNumber} is retained as a string because VCF cardinalities include symbolic values whose interpretation depends on alleles or genotypes; \term{fieldArity} additionally links the declaration to one of nine arity individuals, and \term{fieldNumberInteger} carries the fixed-count case, so that a consumer can branch on cardinality without parsing.

Lexical and parsed values serve complementary purposes. \term{fieldValue} can retain a source token or comma-separated value list, while optional integer, decimal and boolean properties expose suitable scalar interpretations. This makes numeric querying possible without forcing every lexical value into a datatype that cannot represent it: a non-finite VCF floating-point token should not be coerced to \term{xsd:decimal}. Optional raw properties retain header values, INFO and FORMAT order, and condensed sample blocks where textual evidence is needed. Missingness is explicit: a standalone missing value is represented as \term{"."\textasciicircum{}\textasciicircum{}vcfc:Null}, whereas a genotype such as \term{./.} retains its internal allele separators, and in a vector the original missing token occupies its sample position. This avoids interpreting missing depth as zero or conflating an unapplied FILTER with PASS. The meaning of a dot remains context-dependent: \term{Number=.} is a cardinality declaration, not a missing observation. Expanded values may still contain lexical arrays; expansion concerns sample and field granularity, not mandatory decomposition of every VCF substructure.

\section{Two complementary genotype representations}
\label{sec:profiles}

\subsection{Expanded access for single-sample analysis}

With \term{vcfc:ExpandedRepresentation}, a call links through \term{hasSampleCall} to a \term{SampleCall} for each represented sample. Each sample call identifies the sample and links through \term{hasFormatValue} to individual \term{FormatFieldValue} resources. An optional \term{forSample} relation connects it to a reusable sample-column identity. Keeping these resources strictly per-sample preserves their meaning for consumers.

For single-sample files, this structure makes genotype-level access a natural graph traversal. Retrieving calls with read depth at least 20 joins a FORMAT value to the definition whose \term{fieldId} is DP and filters its \term{fieldValueInteger}. Individual calls and values can also receive annotations, which a vector payload cannot carry. The profile remains applicable to larger files when this access pattern justifies the additional graph structure; the vocabulary does not prescribe a measured sample-count threshold.

\subsection{Condensed access for multi-sample cohorts}

The condensed profile factors sample-column identity into one \term{SampleSet}, whose \term{VCFSample} members have an exact \term{sampleName} and a one-based \term{sampleIndex}. A genotype-bearing call links through \term{hasCallMatrix} to a \term{CohortCallMatrix}. This matrix identifies the sample set and groups one \term{FormatValueVector} per represented FORMAT key. Each vector declares its field definition, its \term{valueEncoding} and its \term{encodedValues} payload.

The initial encoding, \term{VCFTextVector}, stores tab-separated raw FORMAT values in sample-index order. Every sample occupies one position, including missing values. Commas within a value remain part of that value, preserving arrays such as allelic depths without confusing them with sample boundaries. The call's \term{formatRaw} retains the order in which keys occur, and \term{sampleDataRaw} optionally preserves the original sample blocks. Condensation therefore retains the represented sample values rather than replacing them with cohort summaries.

This explicit vector class makes the change in query access visible. A vector stores one RDF literal payload, not a collection of independently asserted sample values. A consumer must decode the relevant positions, using the sample order and the field declaration, before performing per-sample comparisons; this can be implemented in application logic or in query functions, but ordinary basic graph patterns alone do not expose those cells. Producers should use the declared profile consistently, and document any deliberate redundant materialization of both forms.

\subsection{Illustration and structural consequences}

Consider a synthetic biallelic record at position 100 on contig 1, with REF=A, ALT=G and FORMAT=GT:DP. Its three sample blocks are \term{0/1:42}, \term{0/0:18} and \term{./.:.}. The condensed representation has the following two payloads, where \term{\textbackslash t} denotes a tab:

\begin{lstlisting}
GT: "0/1\t0/0\t./."
DP: "42\t18\t."
\end{lstlisting}

Selecting position two in each vector recovers SAMPLE2's genotype \term{0/0} and depth 18; the third depth remains missing. An equal-content expanded graph instead has three sample-call resources and six FORMAT-value resources, and a single-sample projection has one sample call and two values. The expanded depth pattern described above selects SAMPLE1, and a condensed consumer obtains the same selection after decoding and excluding the missing depth. This is equivalence of the represented cells under the encoding, not identity of the two RDF graphs.

The structural distinction can be stated independently of a storage engine. Assuming $R$ genotype-bearing records each representing the same $S$ samples and $F$ FORMAT keys, full expanded materialization creates $RS$ sample calls and $RSF$ field-value resources, whereas condensed materialization creates $S$ samples, one sample set, $R$ matrices and $RF$ vectors. Excluding shared file, header and record structures, the resource counts are therefore
\[
N_{\mathrm{expanded}}=RS+RSF,\qquad
N_{\mathrm{condensed}}=S+1+R+RF.
\]
The optional reusable sample layer can also be added to the expanded profile. These counts explain the cohort-oriented design but do not imply a reduction in the number of genotype cells: the condensed payload still contains $RSF$ values. Serialized size, indexing overhead and execution time depend on encoding and implementation, and require measurement.

\section{Specification coverage and version-parameterised validation}
\label{sec:coverage}

\subsection{What was counted, and how}

Claims of comprehensive coverage are not checkable, so we state what was counted, how, and what was excluded. We deliberately do not report a percentage of the specification text: a page or section count mixes normative constructs with worked examples and rationale, and gives a figure that cannot be audited. Instead, the repository carries an inventory of the constructs of the VCF 4.5 \emph{logical model}, derived by walking the specification section by section and recording each thing a conforming representation must be able to express. The inventory contains 106 constructs across seven areas: header line types and their attributes, field declaration attributes, the Type and Number value sets, lexical and encoding rules, the eight fixed fields and allele structure, FORMAT and genotype syntax, and the structural-variant, repeat and gVCF machinery.

Each entry records the specification section, the vocabulary terms that carry the construct, a validation artifact as evidence, and one of three verdicts. \emph{Full} means a dedicated queryable term or term set exists, so a consumer can retrieve the construct with a graph pattern. \emph{Partial} means the construct is representable but only through documentation on a term or through a retained raw literal, so retrieving it requires string handling. \emph{Absent} means no carrier exists. The reserved INFO and FORMAT key declarations are counted separately from the 106 constructs, because they are specification content rather than model structure: a vocabulary that can express arbitrary INFO keys covers the construct whether or not it ships the reserved list.

\begin{table}[t]
\caption{Coverage of the VCF 4.5 logical model, recomputed by \term{npm run coverage:report} from \term{tests/vcf45-coverage-inventory.json}. Reserved key declarations are counted separately in the text.}
\label{tab:coverage}
\centering\small
\begin{tabularx}{\linewidth}{@{}Xrrrr@{}}
\hline
Specification area & Constructs & Full & Partial & Absent \\
\hline
File and header structure & 20 & 20 & -- & -- \\
Field declaration attributes & 8 & 8 & -- & -- \\
Type and arity value sets & 14 & 14 & -- & -- \\
Lexical and encoding rules & 9 & 7 & -- & 2 \\
Fixed fields and allele structure & 18 & 15 & 3 & -- \\
FORMAT and genotype & 21 & 20 & 1 & -- \\
Structural variation, repeats, gVCF & 16 & 16 & -- & -- \\
\hline
\textbf{Total} & \textbf{106} & \textbf{100} & \textbf{4} & \textbf{2} \\
\hline
\end{tabularx}
\end{table}

\subsection{Results and residue}

Table~\ref{tab:coverage} gives the result: 100 constructs fully covered ($94.3\%$), four partially and two absent. Crediting partial constructs at one half yields $96.2\%$; we report both figures rather than choosing the flattering one. The vocabulary declares 445 terms---73 classes, 69 object properties, 108 datatype properties, 182 named individuals, seven annotation properties and five datatypes---across five importable modules covering the core model, alleles, genotypes, structural variation and the reserved-key registry. That registry carries all 122 reserved key declarations of VCF 4.5, 51 INFO and 71 FORMAT, each with the Number, Type and Description given in the specification; it is generated from the specification's own source rather than transcribed, and the generated file records the SHA-256 of the input it was derived from.

The residue is as informative as the total. The two absent constructs are file-level byte state: whether the source was UTF-8 with or without a byte order mark, and whether lines ended with LF or CR+LF. Both are properties of the octet stream rather than of the logical model, and a carrier for them would be a claim about a file the RDF graph no longer has. The four partial constructs are recorded in prose on the relevant term rather than as separate resources: the angle-bracketed assembly-contig form of CHROM, the telomeric POS values 0 and $N+1$, the REF padding-base convention, and the distinction between a single missing token and a per-item missing list. Each is expressible today by reading a documented term comment; none is reachable by a graph pattern, which is why they are not counted as full.

Three areas are excluded by design rather than by omission, and the inventory records the reason for each. The BCF~2.2 byte layout is out of scope because the vocabulary targets the VCF logical model and not the binary serialization, even though the two are specified in one document. Byte-for-byte source reconstruction is out of scope because the raw carriers preserve logical content rather than file bytes. A VCF~4.0 conformance profile is not published, and the validator reports that version as unsupported rather than silently validating it against later rules.

\subsection{Reproducing the figures}

Every figure in this section is produced by \term{npm run coverage:report}. The script reads the five ontology modules, the SHACL profiles, the per-version reserved-key registries and the inventory; it writes a machine-readable \term{tests/coverage-report.json} containing the per-area counts, the term census, the external-vocabulary breakdown and the validation summary; and it exits non-zero if the inventory cites a vocabulary term that the ontology no longer declares. The verdicts remain editorial judgements and we do not claim otherwise, but they are individually attributable to a construct, a section and an artifact, and neither the arithmetic nor the term references can silently drift from the vocabulary. The command runs inside the repository's validation pipeline, so a change that invalidates a published figure fails the build rather than surviving to the next paper.

\subsection{One vocabulary, five conformance profiles}

VCF versions differ less in what they can express than in what they permit. Between 4.1 and 4.5 the record structure, the header grammar and the genotype block are stable; what changes is which reserved keys exist, which Number codes are legal, and how wide the structural-variant aggregates are. Modelling each version as its own vocabulary would duplicate the stable majority, fragment the term space, and force consumers to translate between near-identical vocabularies to compare files that differ only in vintage. We therefore treat version differences as constraint differences: one vocabulary describes the shared model, and validation is parameterised by version.

Five SHACL overlays, \term{vcf-4.1.shacl.ttl} through \term{vcf-4.5.shacl.ttl}, carry the per-version rules. Their constraints are conditioned on the owning file's \term{fileFormat}, so all five may be loaded together and each applies only to the files it governs; the optional \term{VCF41File}\ldots\term{VCF45File} classes additionally check that a file's declared version matches its intended profile. A representative difference: CIPOS carries two values in 4.1--4.3 but two per alternate allele in 4.4--4.5, so the same aggregate-width rule is expressed twice, with different arithmetic and different version guards, and a 4.2 file is not judged by a 4.5 rule. Number codes behave the same way: the local-allele codes LA, LR and LG and the base-modification code M are gated to the versions that define them.

The reserved-key registries are versioned as well, and their sizes track the specification's growth---40 declarations for 4.1 and 4.2, 67 for 4.3, 78 for 4.4 and 122 for 4.5. Each snapshot records the official source URL and a SHA-256 of the specification text it was derived from, so a reader can confirm which revision a rule reflects. For 4.1 and 4.2, only explicit declarations are captured; prose-only INFO fields whose exact format the specification leaves to producers are not assigned invented Number or Type requirements, because doing so would manufacture conformance failures for files that were always valid.

The overlays sit on top of three version-independent profiles: a portable structural profile of 53 node shapes and 115 property shapes, a SPARQL profile of 16 cross-resource ordering and uniqueness rules, and a consistency profile of 21 rules relating raw tokens to their parsed counterparts. This layering keeps the portable profile usable by validators without SPARQL support, keeps cross-resource rules optional rather than mandatory, and makes a version upgrade a matter of adding an overlay rather than revising the model.

\section{Implementation and validation}

The vocabulary is already used as the target model of VCF-RDFizer, an open-source VCF-to-RDF converter distributed under the MIT license~\cite{vcfrConverter2026}. Its default RML mappings construct file, header, record and call resources in the vocabulary's namespace. Complementary streaming emitters produce per-sample \term{SampleCall} and \term{FormatFieldValue} resources, or sample-ordered \term{CohortCallMatrix} and \term{FormatValueVector} resources, selected through the converter's \term{--sample-representation} option. Configurable mappings, documented execution paths and regression tests for the two genotype structures provide concrete evidence that the model can guide an implemented conversion workflow. Wider adoption, complete conformance and comparative performance remain separate evaluation objectives; in particular, the converter's optional RDF compression stages operate on the resulting graph and should be evaluated separately from the vocabulary's choice of genotype granularity.

Version 2.0.0 provides OWL/Turtle definitions, SHACL shapes, example graphs, illustrative mapping patterns and documentation generated from the vocabulary source~\cite{vcfc200}. The synthetic example accompanying this manuscript provides both profiles and an application-level check that recovers all six FORMAT cells and the same depth-based sample selection. This is a small representation check, not a performance or general conformance study.

The constraint layer is executed rather than merely published~\cite{shacl2017}. The suite validates 20 fixtures---the maintained examples, both paper graphs, the per-version examples and a set of negative probes---with zero violations and zero warnings, and the regression tests exercise each constraint in isolation, including shared shape dependencies, so that an unrelated violation cannot mask a broken rule. Decoded checks that SHACL cannot express portably, such as genotype-likelihood cardinalities at arbitrary ploidy and Number=M site counts, run in a complementary Python layer and are reported separately, so it remains visible which layer established which guarantee. Structural validation stays distinct from validation of encoded VCF contents: the shapes check that a graph is well-formed against the specification's rules, not that it faithfully reproduces a particular file.

The model's scope still bounds fidelity claims. Alternate-allele order, which repeated RDF assertions alone cannot preserve, is carried by explicit \term{AltAllele} resources with a zero-based \term{alleleIndex} to which the Number=A, R and G families bind their values; but raw fields remain optional, so one valid graph may retain less source text than another. File-based identifiers provide local traceability and still require a publication policy for stable global identification. Coverage as reported in Section~\ref{sec:coverage} is coverage of the specification's constructs by vocabulary terms and validation rules---it is neither a conformance certification against a VCF test corpus nor a demonstration of interoperability, and universal byte-for-byte reconstruction remains explicitly out of scope. The condensed profile remains an additive extension: existing per-sample terms keep their meanings, and declaring the encoding on each vector separates the representation contract from a particular payload syntax, so further encodings can be introduced provided their ordering, missingness and decoding rules are specified.

The vocabulary nevertheless fixes enough of the representation contract to support three concrete investigations that hold the target constant while varying one factor. First, a fidelity suite can compare reconstructed values for phased and partially missing genotypes, variable ploidy, multiallelic and structural variants, and heterogeneous header declarations, distinguishing semantic value recovery from exact source serialization. Second, equal-content conversion experiments can vary records, samples, FORMAT fields and missingness, measuring triples, serialized and compressed bytes, peak memory, and conversion and loading time against VCF/BCF baselines. Third, query experiments can compare direct expanded access with vector decoding and selective expansion, checking identical answer sets before assessing execution time. External alignment experiments can then assess reference-aware variant joins and federated use separately from conversion fidelity.

\section{Conclusion}

The VCF Core Vocabulary establishes a versioned foundation for representing VCF observations together with their interpretation context. It reuses published vocabularies where they already model the domain---FALDO directly for positions and strandedness, SO, VRS, GENO, ChEBI and HERO as alignment targets---and reserves minted terms for VCF syntax itself. It covers 100 of the 106 constructs of the VCF 4.5 logical model on a published inventory that is recomputed by the repository's own validation pipeline, with its remaining gaps and its excluded scope stated explicitly rather than left to the reader. It treats VCF versions as a validation parameter, so that one model serves 4.1 through 4.5 through five SHACL conformance profiles rather than five parallel vocabularies. And its two genotype profiles let a producer choose direct access to individual sample calls or an explicit cohort structure that does not materialize every sample/field intersection. Use in the VCF-RDFizer converter provides a concrete implementation setting for comparing fidelity and computational behavior against this shared semantic foundation.

\begin{acknowledgments}
Project funding provided from the Research Foundation -- Flanders (FWO) (SB Fellowship 1S27825N).
Ruben Taelman is a postdoctoral fellow of the Research Foundation -- Flanders (FWO) (1202124N).
% KNoWS-IDLab project attribution template. Add the project name and grant
% agreement number if this work is booked to a named KNoWS-IDLab project rather
% than the first author's PhD; see ACKNOWLEDGEMENTS.md in the repository.
%   This work was carried out at KNoWS, IDLab, Ghent University -- imec, within
%   <KNoWS-IDLab project name> <grant agreement number>.
\end{acknowledgments}

\section*{Declaration on Generative AI}
OpenAI Codex and Claude assisted with literature discovery and reference checking, outline development, text drafting, preparation of the schematic and synthetic example, the 2.0.0 rename, the VCF 4.5 construct inventory and its reproducible measurement scripts, and LaTeX verification. This preparatory draft requires substantive review and revision by the human authors before submission; in particular, the coverage verdicts recorded in the inventory are editorial judgements that the authors must confirm.

% Finalize this declaration after author revision; do not claim completed human
% review prematurely. See author-notes.md for the current CEUR-WS restrictions.

% Main matter currently runs to 9 pages against the 5-page CEUR limit; the
% authors chose completeness over length for this draft. Sections 2.2, 4 and 5
% are the candidates for compression before submission.
% References are excluded from the five-page main-matter limit.
\clearpage
\bibliography{sources}
\end{document}
